\section{ТЕСТИРОВАНИЕ И АНАЛИЗ РЕЗУЛЬТАТОВ}

\subsection{Организация процесса тестирования}
Для обеспечения высокого качества разработанного программного обеспечения и веб-ресурса был проведен комплексный цикл тестирования. Тестирование осуществлялось на нескольких уровнях: функциональное тестирование компонентов, кроссбраузерное тестирование верстки и альфа-тестирование на фокус-группе.

Функциональное тестирование Telegram-бота (Unit-подобное ручное тестирование) проводилось по методу «черного ящика». Были составлены тест-кейсы для проверки всех возможных путей прохождения диалогов (Happy Path и альтернативные сценарии). Особое внимание уделялось тестированию машины состояний (FSM). Проверялась корректность обработки непредвиденного пользовательского ввода: например, когда пользователь отправляет текстовое сообщение в тот момент, когда бот ожидает нажатия инлайн-кнопки (Inline Button), или вводит нечисловые данные там, где ожидается число. Бот успешно справился с обработкой исключений (Exceptions), возвращая пользователю понятные сообщения об ошибке и не прерывая сессию.

Кроссбраузерное тестирование веб-сайта проводилось с использованием инструментов разработчика (DevTools) в браузерах Google Chrome, Mozilla Firefox и Safari. Проверялась корректность отображения эффекта Glassmorphism, так как свойство \texttt{backdrop-filter} исторически имело проблемы с поддержкой в старых версиях браузеров. Для обеспечения обратной совместимости (graceful degradation) были предусмотрены резервные (fallback) сплошные полупрозрачные фоны.

\subsection{Результаты пользовательского тестирования (Альфа-тест)}
После завершения внутреннего тестирования проект был представлен небольшой фокус-группе, состоящей из 15 студентов первого курса. Целью данного этапа было оценить удобство пользовательского интерфейса (UX/UI) сайта и степень вовлеченности аудитории при взаимодействии с Telegram-ботом.

Участникам было предложено:
\begin{enumerate}
    \item Ознакомиться с материалами на веб-сайте.
    \item Перейти в Telegram и пройти тест на цифровую гигиену.
    \item Попытаться намеренно «сломать» логику бота (stress testing).
\end{enumerate}

\textbf{Анализ результатов тестирования показал следующее:}\\
1. 100\% участников отметили привлекательный и современный дизайн сайта. Темная тема оформления была оценена крайне положительно, особенно при чтении материалов с экранов мобильных телефонов.\\
2. 85\% пользователей полностью прошли тест на цифровую гигиену и дошли до финального сообщения с результатами.\\
3. По результатам тестирования цифровой гигиены выяснилось, что более 60\% участников фокус-группы находятся в «зоне риска» — они проводят за экранами более 8 часов в сутки и не используют фильтры синего света перед сном. Этот факт эмпирически подтвердил высокую актуальность выбранной темы проекта.\\
4. В ходе стресс-тестирования бота не было зафиксировано ни одного критического сбоя (Crash) интерпретатора Python, что подтверждает надежность заложенной архитектуры.

\subsection{Развертывание (Deployment) и эксплуатация}
Итоговая версия исходного кода проекта, включая все HTML/CSS файлы сайта и Python-скрипты Telegram-бота, была зафиксирована в системе контроля версий Git. Код был отправлен (push) в удаленный репозиторий на платформе GitHub.

Для запуска Telegram-бота в режиме промышленной эксплуатации (production) достаточно развернуть скрипт на любом облачном VPS (Virtual Private Server) под управлением операционной системы Linux. Благодаря наличию файла \texttt{requirements.txt}, установка зависимостей (\texttt{pip install -r requirements.txt}) и запуск бота занимает менее двух минут. Статический веб-сайт полностью готов к автоматическому развертыванию через сервис GitHub Pages без дополнительных настроек.
